\documentclass[authoryear,journal,12pt]{elsarticle}

\usepackage{graphicx,natbib,amssymb,lineno}
\usepackage{amsmath}
\usepackage{amsbsy,amsfonts}
\usepackage{bm}
\usepackage{epsfig}
\usepackage{soul}
\usepackage[normalem]{ulem}
\usepackage{booktabs}
\usepackage{array,multirow}
\usepackage[percent]{overpic}
\usepackage{epstopdf}
\usepackage{hyperref}

\begin{document}

\begin{frontmatter}

\title{Active Perception of Object Shape Using a Sensorless Dexterous Robotic Hand}

\author[sch]{G.~Sai~Pranav}
\author[iitgn]{Prof.~Harish J.~Palanthandalam-Madapusi}

\begin{abstract}

Robots commonly rely on cameras, depth sensors, or dedicated tactile sensors
to identify and understand objects. In this project, we investigate whether
a dexterous robotic hand can identify an object's shape and estimate its
dimensions using proprioceptive information obtained through physical
interaction, without relying on external visual or dedicated tactile
sensing.

We use a LEAP Hand to interact with two geometric objects: a cube with a
side length of 7.5~cm and a sphere with a diameter of 7.5~cm. The hand
performs controlled exploratory movements while we record proprioceptive
information such as joint states, motor positions, finger trajectories,
and, where applicable, motor currents. These observations provide indirect
information about the object's geometry through the way the fingers move
and interact with its surface.

We follow an active perception approach in which manipulation is treated as
a source of information. Rather than making a decision from a single grasp,
the hand explores the object through successive interactions and uses the
resulting observations to construct an initial internal representation. As
manipulation continues, new contacts generated through movements such as
rotation, sliding, pivoting, or regrasping provide additional information
that can be used to refine this representation. We then identify the object
as a cube or a sphere and estimate its corresponding physical dimension.

\end{abstract}


\end{frontmatter}

\textbf{Keywords:} Active perception, dexterous manipulation, proprioception,
robotic hand, object shape recognition, LEAP Hand, proprioceptive perception.

\clearpage

\renewcommand{\baselinestretch}{1.0}
\section{Introduction}
\label{sec:intro}

Understanding an unfamiliar object is an important capability for a robot
that needs to interact with its surroundings. Humans can often infer an
object's approximate shape and size through physical interaction by moving
their fingers, changing their grasp, and observing how the object responds,
even when visual information is unavailable. This raises the possibility
that a robotic hand could similarly obtain information about an object
through its own movements and interactions.

Dexterous robotic hands provide a suitable platform for investigating this
idea. Their multiple fingers and joints generate proprioceptive information
about the configuration and movement of the hand. When the fingers interact
with an object, these measurements can indirectly reflect characteristics
of its geometry.

Previous studies have demonstrated the potential of proprioception for
object recognition. V'asquez and Perdereau developed a proprioceptive
shape signature based on joint angles to recognize primitive object shapes
and estimate object pose. Their representation was designed to be invariant
to object size and position \citep{VasquezPerdereau2017}. Sintov and Meir
subsequently demonstrated object recognition using the kinematics of a
multi-fingered robotic hand without relying on visual or tactile sensing,
including recognition of broader geometric categories
\citep{SintovMeir2023}.

These studies establish that proprioceptive information can provide useful
information about object geometry. However, their primary focus is on
recognizing object shapes or categories, rather than estimating the
physical dimensions of the object as an output. The present project builds
on this direction by investigating whether proprioceptive interaction can
be used for both shape identification and dimensional estimation.

In the present study, we investigate whether a dexterous robotic hand can
identify an object as a cube or sphere and estimate its physical dimension
using proprioceptive information obtained through active physical
interaction, without external visual or dedicated tactile sensing. We use
a LEAP Hand and two geometric objects: a cube with a side length of 7.5~cm
and a sphere with a diameter of 7.5~cm. Their known dimensions provide ground
truth against which the estimated dimensions can be evaluated.


\section{Data and Methods}
\label{sec:methods}

\subsection{Experimental Setup}

The experiments were conducted using a LEAP Hand, an open-source
dexterous robotic hand with four fingers and 16 actuated joints. The
index, middle, and ring fingers each have four joints, while the thumb has
four joints comprising base rotation and flexion-related movements. The
hand provides proprioceptive information about its joint configuration
during physical interaction with an object.

Two geometric objects were used in the experiments: a cube with a side
length of 7.5~cm and a sphere with a diameter of 7.5~cm. The known
dimensions of these objects provide ground-truth values for evaluating
the estimated dimensions.

The hand was operated using the Ex-ARM Python control framework developed
for the LEAP Hand. The framework provides an interface for hardware
control and includes a kinematic model for calculating fingertip
positions from the joint configuration. The hardware interface provides
access to joint positions, velocities, and motor currents, while the
kinematic model provides forward kinematics for the four fingertips.

\subsection{Proprioceptive Data}
\label{sec:proprioceptive_data}

The perception process relies on information obtained from the robotic
hand itself rather than from external cameras, depth sensors, or dedicated
tactile sensors. During physical interaction, the configuration of the
fingers changes according to the geometry of the contacted surface. These
changes provide indirect information about the shape and dimensions of
the object.

The primary proprioceptive measurements consist of joint states obtained
from the LEAP Hand. Motor-current measurements are also used to identify
changes associated with physical contact. Since motor current is not a
dedicated tactile measurement, it is treated as a contact proxy rather
than as a direct measurement of contact force.

The recorded joint positions are supplied to the LEAP Hand forward
kinematic model to determine the corresponding fingertip positions in the
hand coordinate frame. The resulting fingertip locations provide spatial
observations of the object surface during interaction.

\subsection{Active Physical Exploration}
\label{sec:active_exploration}

The object is explored through a sequence of physical interactions rather
than through a single static grasp. The fingers are moved towards the
object and contact is identified from changes in the proprioceptive state.
For each detected contact, the corresponding joint configuration is used
to determine the fingertip position through forward kinematics.

The hand and object are maintained in a fixed relative configuration
during individual probing operations, and the hand is released between
successive contacts. This allows contact observations to be collected
from different probing configurations while maintaining a common
reference frame.

Each new contact provides an additional observation of the object's
surface. Consequently, the geometric representation is progressively
constructed from multiple contact locations. Movements that change the
contact configuration can expose different portions of the object and
provide additional information for perception.

This approach treats manipulation as part of the perception process.
Instead of separating perception from manipulation, the physical
interaction itself provides the observations from which the object's
geometry can be inferred.

\subsection{Object Shape Identification}
\label{sec:shape_identification}

The collected proprioceptive observations are used to distinguish between
the two object geometries. The fingertip contact locations obtained from
the hand's joint states provide a spatial representation of the contacted
surface.

For the sphere, contact locations are expected to follow the curved
surface of the object, whereas contacts with the cube are associated with
its planar faces and edges. The spatial distribution of the contact
observations is therefore used to distinguish the two geometric classes.

The Ex-ARM implementation includes an experimental proprioceptive
shape-recognition procedure based on LEAP joint states and motor-current
changes for distinguishing sphere-like and box-like objects. Contact
samples are generated from the measured joint positions using the LEAP
Hand forward-kinematics model.

\subsection{Dimensional Estimation}
\label{sec:dimension_estimation}

Following shape identification, the proprioceptive contact observations
are used to estimate the corresponding physical dimension of the object.
For the cube, the characteristic dimension is its side length, while for
the sphere it is its diameter.

The estimated dimension is obtained from the spatial information provided
by the successive contact observations and is compared with the known
ground-truth dimension of 7.5~cm. This provides a direct measure of how
well proprioceptive interaction can recover the physical scale of the
object without external visual or dedicated tactile sensing.

The complete implementation used in this study is publicly available in
the Ex-ARM GitHub repository \citep{ExARM}.

\section{Results}
\label{sec:results}

\clearpage

\bibliographystyle{elsart-harv}
\bibliography{manuscript}

\end{document}